\centerline{\Huge Условия и решения задач.}
\centerline{\bf \Huge 4 класс}

\problem{4.1} В подземном царстве одну золотую монету можно обменять на четыре серебряных, а десять серебряных монет на один алмаз. Сколько алмазов можно получить, имея пять золотых монет? \par

\textbf{Решение.} Пять золотых монет можно обменять на $5\cdot4 = 20$ серебряных монет. В свою очередь эти 20 серебряных можно обменять на $20:10 = 2$ алмаза.  \par

\textbf{Критерии.} \\
\textit{7 баллов} $-$ Полностью верное решение.\\
\textit{3 балла} $-$ Получил, что 5 золотых монет это 20 серебряных.\\
\textit{0-2 балла} $-$ Есть какие-то продвижения, которые могут привести к решению. \\ 
\textit{0 баллов} $-$ Решение отсутствует. Решение неверное, продвижения отсутствуют. \\


\problem{4.2} Володя взял из большого мешка половину шоколадок и переложил их в четыре маленьких мешочка. Однако Володя передумал, достал по 7 шоколадок из каждого маленького мешочка и положил обратно в большой мешок. На сколько в большом мешке стало больше шоколадок, чем в маленьких мешочках? \par

\textbf{Решение.} После того как Володя передумал, в большом мешке помимо той половины, что в ней была, оказалось ещё $7\cdot4$ шоколадок. То есть «половина» плюс $7\cdot4$. А в 4 маленьких мешках суммарно была половина, но после того как Володя передумал, осталось «половина» минус $7\cdot4$. Чтобы узнать на сколько больше шоколадок в большом мешке, чем в маленьких нужно от первого отнять второе. То есть  («половина» $+$ $7\cdot4$) $-$ («половина» $-$ $7\cdot4$) $=$ «половина» $+$ $7\cdot4$ $-$ «половина» $+$ $7\cdot4$. Видно, что +«половина» и -«половина» сокращаются и остается 56. Ответ: на 56 шоколадок больше. \\

\textbf{Критерии.} \\
\textit{7 баллов} $-$ Полностью верное решение.\\
\textit{4 балла} $-$ Правильно выражено, сколько в маленьких мешках \textbf{И} в большом мешке шоколадок через половину большого.\\
\textit{2 балла} $-$ Правильно выражено, сколько в маленьких мешках \textbf{ИЛИ} в большом мешке шоколадок через половину большого.\\
\textit{0-2 балла} $-$ Есть какие-то продвижения, которые могут привести к решению. \\ 
\textit{0 баллов} $-$ Решение отсутствует. Решение неверное, продвижения отсутствуют. \\ 

\problem{4.3} В Калмыкии между любыми двумя населенными пунктами есть ровно одна дорога. Докажите, что найдутся два населенных пункта, из которых исходит одинаковое количество дорог. \par

\textbf{Решение.} Пусть в Калмыкии $n$ населённых пунктов. Из каждого населенного пункта может исходить либо 0, либо 1, либо 2, $\dotsc$, либо $n-1$ дорога. Заметим, что не может одновременно существовать два населённых пункта из которого исходит 0 и из которого исходит $n-1$ дорога. Так как первый не будет соединен ни с кем, а второй должен быть соединен со всеми, а значит и с первым тоже. То есть из этих двух вариантов может существовать только один. Тогда у каждого населенного пункта $n-1$ вариант, сколько из него будет исходить дорог. Населенных пунктов у нас $n$, а количество вариантов исходящих из них дорог $n-1$. Значит по принципу Дирихле есть два пункта, их которых исходит одинаковое число дорог. \\ 

\textbf{Критерии.} \\
\textit{7 баллов} $-$ Полностью верное решение. \\
\textit{4 балла} $-$ Доказано, что не может существовать пунктов с исходящими 0 и $n-1$ дорогами одновременно.\\ 
\textit{2 балла} $-$ Написано, что из каждого населенного пункта может исходидить либо 0, либо 1, $\dotsc$, либо $n-1$ дорога. \\
\textit{0-2 балла} $-$ Есть какие-то продвижения, которые могут привести к решению. \\ 
\textit{0 баллов} $-$ Решение отсутствует. Решение неверное, продвижения отсутствуют. \\ 

\problem{4.4} У Аркадия есть клетчатая доска $23\times23$, в каждой клетке которой сидит по кузнечику. В какой-то момент времени каждый кузнечик перепрыгнул на соседнюю по стороне клетку. Докажите, что после этого найдется хотя бы одна пустая клетка.

\textbf{Решение.} Раскрасим доску в шахматном порядке в красный и синий цвета. Заметим, что 23 это нечётное число. Следовательно $23\cdot23$ тоже нечётное. Пусть не нарушая общности красных клеток больше чем синих(на одну). Заметим, что кузнечик сидящий на красной клетке, после прыжка будет сидеть на синей и наоборот. Стало быть синие кузнечики пересели на все красные клетки. Но так как красных клеток на одну больше чем синих, то даже если синие кузнечики уселись на разные места, останется хотя бы одно свободное место. Что и требовалось доказать. \\

\textbf{Критерии.} \\
\textit{7 баллов} $-$ Полностью верное решение.\\ 
\textit{4 балла} $-$ Доказано, что кузнечник меняет цвет клетки при прыжке.\\ 
\textit{2 балла} $-$ Правильно выбранная раскраска.\\
\textit{0-2 балла} $-$ Есть какие-то продвижения, которые могут привести к решению.\\
\textit{0 баллов} $-$ Решение отсутствует. Решение неверное, продвижения отсутствуют. \\ 

\problem{4.5} На острове Парадиз проживают $N$ островитян: эльдийцы, которые всегда говорят правду и марлийцы, которые всегда лгут. Путешственник задал им вопрос: «Кого на острове больше $-$ эльдийцев или марлийцев?» \\
Первый житель острова ответил «Марлийцев больше» \\ 
Второй $-$ «Эльдийцев не меньше» \\
Третий «Марлийцев больше» \\
Четвертый «Эльдийцев не меньше» \\
… \\
$(N-1)-$ый ответил «Эльдийцев не меньше» \\
$N-$ый «Марлийцев больше». \\
Может ли $N$ равняться $2023$? \par

\textbf{Решение.} Заметим, что высказывание «Марлийцев больше» противоположно высказыванию «Эльдийцев не меньше». Стало быть одно из них правда, а другое ложь. Следовательно среди отвечавших марлийцы и эльдийцы чередовались. Предположим, что $N = 2023$. Тогда высказываний «Марлийцев больше» было произнесено $1012$ раз, а «Эльдийцев не меньше» $1011$ раз. Пусть «Марлийцев больше» говорили эльдийцы, которые всегда говорят правду. Но тогда эльдийцев $1012$ и они соврали, но они всегда говорят правду. Противоречие. Пусть «Марлийцев больше» говорили марлийцы, которые всегда лгут. Тогда марлийцев $1012$ и они сказали правду, но они всегда лгут. Противоречие. Следовательно $N$ не может равняться $2023$. \\

\textbf{Критерии.} \\
\textit{7 баллов} $-$ Полностью верное решение.\\ 
\textit{4 балла} $-$ Понимание того, что высказывания чередуются среди марлийцев и эльдийцев.\\
\textit{3 балла} $-$ Доказано, что высказывания «Марлийцев больше» и «Эльдийцев не меньше» противоречивы.\\
\textit{0-2 балла} $-$ Есть какие-то продвижения, которые могут привести к решению. \\
\textit{0 баллов} $-$ Решение отсутствует. Решение неверное, продвижения.\\
